Determinar la massa i les posicions dels cossos celestes va ser, sens dubte, un gran repte per als primers astrònoms. Gràcies a les valuoses dades sobre les posicions dels astres que Tycho Brahe va recollir al llarg de la seva vida, Johannes Kepler va poder formular les seves famoses tres lleis.

\figura[0.4]{fig1.png}
{\centering Kepler (esquerra) i Brahe (dreta)\par}

\begin{apartat}{1}
Deduïu la tercera llei de Kepler a partir de la segona llei de Newton i de la llei de la gravitació universal, suposant que els planetes descriuen moviments circulars uniformes al voltant del Sol.
\end{apartat}

\begin{apartat}{1}
Determineu la massa del Sol emprant les dades que necessiteu de la taula següent:

\smallskip
\begin{tabular}{|l|c|c|}
\hline
\textit{Planeta} & \textit{Radi de l'òrbita} (\textit{m}) & \textit{Període} (\textit{anys})\\
\hline
Mercuri & $57,90\times 10^{9}$ & 0,2408\\
\hline
Venus & $108,2\times 10^{9}$ & 0,6152\\
\hline
Terra & $149,6\times 10^{9}$ & 1,000\\
\hline
Mart & $228,0\times 10^{9}$ & 1,881\\
\hline
\end{tabular}
\end{apartat}

\blocetiquetat{Dada:}{$G = 6,67\times 10^{-11}\ \mathrm{N\,m^2\,kg^{-2}}$.}
